\section{The complete prime-word graph and the fate of every packet class}
\label{sec:word-graph-completion}

This calculation uses the literal source family of the received
\emph{Constructive two-prime packets: exact theta lifts and bounded
moment dynamics}, results SZ-20260923-901--903. Both of its complete
TeX sources, \path{NOTE.tex} and \path{GENERATOR.tex}, were read.
Their local source identities and original human reading coverage
are retained in the private source ledger \cite{ConstructiveTwoPrime}. The mathematical context
for the original theta quotient is Meyer \cite{Meyer}; the entire
primitive section used here is proved in PL1--27 above. The present
chapter computes the completions and their exact maps, using the
new word family without changing its source metric.

\subsection{The original source and its full finite observation}
Retain the finite reflected packet $Z$, all multiplicities
$m_\rho$, the original coordinates $t_\rho=s-\rho$, and
$h(s)=c_h\prod_\rho(s-\rho)^{m_\rho}$. Put
$d=\sum_\rho m_\rho$ and $A=\prod_\rho\C[t_\rho]/(t_\rho^{m_\rho})$.
The primitive map $S:A\to L^2(\R_{>0},dx)$ is exactly
$S u=\operatorname{Prim}_h(u)$ of PL; its Mellin transform is
\[
 \mathcal M(Su)(s)=a(s)\frac{g(s)}{h(s)}R_u(s),\qquad
 a(s)=s(s-1),\quad g=2\xi,\quad
 R_u=\operatorname{rem}_h([a(g/h)]^{-1}u).
 \tag{GC1}\label{gc:section}
\]
It has full selected jets $u$ and vanishes at every unselected
zero to its full order. Consequently $S$ is injective. Its
positive Gram $G=S^*S$ retains the original measure $dx$ and all
original local units; it is not the reflected Weil form.
The latter's exact connecting formula is GC11 below.

Let $H_c=a_2+c b_2+c^2a_3+c^3b_3$, where
$\lambda_p=p^{s-1/2}$, $a_p=(\lambda_p+\lambda_p^{-1})/2$ and
$b_p=(\lambda_p-\lambda_p^{-1})/(2i)$. Choose the finite selector
$c$ of CG3--5 in the received proof. To recall why it exists,
the original prime characters for 2 and 3 separate every pair of
distinct roots: equality would imply an impossible rational
relation between $\log2$ and $\log3$. Each value collision or
vanishing derivative is therefore the zero of a nonzero polynomial
in $c$ of degree at most three. Excluding the derivatives only at
multiple roots leaves at most $3(\binom r2+r_2)$ bad integers,
where $r=|Z|$ and $r_2$ counts the multiple roots. One of the
integers from zero to that number avoids them all.
Thus $z_\rho=H_c(\rho)$ are distinct, and the local linear term
of $H_c-z_\rho$ is nonzero at each nonsimple factor. Substitution
there is triangular with invertible diagonal; interpolation at
the distinct nodes then gives the unital surjection
\[
 E:\C[X]\twoheadrightarrow A,\qquad P\longmapsto P(H_c),
 \qquad \ker E=(\nu),\quad
 \nu(X)=\prod_\rho(X-z_\rho)^{m_\rho}\in\R[X].
 \tag{GC2}\label{gc:observation}
\]
For completeness, the interpolation assertion follows by solving
the finite inverse of each local substitution coefficient in
increasing degree, then using the relatively prime powers
$(X-z_\rho)^{m_\rho}$. Their product divides exactly the
polynomials whose full substituted jets vanish. The resulting
map is onto and its kernel is the displayed degree-$d$ ideal.
The conjugation $H_c^\dagger=H_c$ makes the nodes conjugate in
pairs with equal multiplicities, proving the real coefficients
of $\nu$. No nilpotent order has been removed.

Write $\mathsf H_c$ for the actual four-dilation operator:
\[
 \mathsf H_c=\tfrac{1-ic}{2}V_2+\tfrac{1+ic}{2}V_{1/2}
       +\tfrac{c^2-ic^3}{2}V_3+\tfrac{c^2+ic^3}{2}V_{1/3},
 \qquad (V_qF)(x)=q^{-1/2}F(x/q).
 \tag{GC3}\label{gc:operator}
\]
Each $V_q$ is unitary for $dx$, with adjoint $V_{q^{-1}}$.
Therefore $\mathsf H_c$ is bounded and self-adjoint. Its real
Mellin-line multiplier and bound are
\[
 \begin{aligned}
 x_c(t)&=\cos(t\log2)+c\sin(t\log2)
                +c^2\cos(t\log3)+c^3\sin(t\log3),\\
 |x_c(t)|&\le K_c=(1+c^2)^{3/2}.
 \end{aligned}
 \tag{GC4}\label{gc:symbol}
\]
Put $F_*=S1$ and $\Phi(P)=P(\mathsf H_c)F_*$. Mellin
Plancherel, with its original factor, gives the finite positive
measure
\[
 \int f(x)\,d\mu(x)=\frac1{2\pi}\int_\R
       f(x_c(t))|\mathcal MF_*(1/2+it)|^2\,dt,
 \qquad \|\Phi(P)\|^2=\int|P|^2\,d\mu.
 \tag{GC5}\label{gc:measure}
\]
Its mass is $\|F_*\|^2$, and its support lies in $[-K_c,K_c]$.
The original source relation from PL and the received proof is
$\Phi(P)=S(EP)+\Theta\mathcal D_P(1)$ with its explicitly
integrated source $\mathcal D_P$. In particular
$\Phi(\nu Q)=\Theta\mathcal D_{\nu Q}(1)$ is an actual theta
relation, not merely a vanished finite observation. Conversely
every theta relation has zero selected jets, so the kernel in
GC2 is exactly the relation ideal in this word family.

\subsection{Every native minimum tends to zero}
The measure $\mu$ has no atoms. Indeed $x_c$ is a nonconstant
real-analytic function: its four distinct exponential frequencies
are linearly independent, and the coefficient of the frequency
$\log2$ is nonzero. Their linear independence follows by
differentiating a vanishing linear combination zero through
order three; the resulting Vandermonde determinant is nonzero.
Every level set of a nonconstant real-analytic function of one
variable is discrete and therefore countable. Its inverse
image under GC5 has Lebesgue measure zero, proving the assertion.
Thus $\mu(\nu^{-1}(0))=0$.

Polynomials are dense in $\mathcal H=L^2(\mu)$.
Here are the approximation details used in this statement.
For a finite Borel measure on a compact interval, continuous
functions are dense in $L^2$: approximate the indicator of an
open set $O$ by $\min(1,n\operatorname{dist}(x,O^c))$ and use
dominated convergence. Complements, finite intersections (take
products of bounded approximants) and increasing countable
unions show that indicators of all Borel sets are in the same
closure; simple-function approximation finishes the claim.
After an affine map to $[0,1]$, the Bernstein polynomials of a
continuous $f$ approximate it uniformly. Indeed their error is
bounded by its modulus of continuity at $\delta$, plus
$2\|f\|_\infty/(4n\delta^2)$, by the mean $x$ and variance
$x(1-x)/n$ of the binomial weights. First choose $\delta$ small
and then $n$ large. This proves the required polynomial density.

Multiplication by $\nu$ on $\mathcal H$ is bounded and
self-adjoint, with zero kernel, because $\nu$ is real on the
interval and vanishes on a set of zero measure. Its range is
dense: the orthogonal complement of the range is its adjoint's
kernel. Continuity of multiplication and polynomial density
therefore imply
\[
 \overline{\nu\C[X]}^{\,L^2(\mu)}=\mathcal H.
 \tag{GC6}\label{gc:relation-closure}
\]

For $n\ge d-1$, let $G_n$ be the attained minimum Gram on $A$
in its original ascending jet coordinates:
\[
 \langle G_nu,u\rangle
 =\min\{\|P\|_{L^2(\mu)}^2:\deg P\le n,\ EP=u\}.
 \tag{GC7}\label{gc:minimum}
\]
Each is positive definite by polynomial injectivity in $L^2(\mu)$;
that injectivity follows because $\mu$ has infinitely many
support points whereas a nonzero polynomial has finitely many
zeros. The feasible space is a nonempty finite-dimensional affine
space, so the minimum is attained. The Grams decrease in the
positive-semidefinite order. For a fixed representative $P_u$,
GC6 approximates $-P_u$ by $\nu Q_j$ in $L^2(\mu)$.
Then $P_u+\nu Q_j$ retains exactly $u$ and its norm tends to zero.
Consequently $G_n\downarrow0$ as $n\to\infty$.
This convergence is in operator norm in the finite-dimensional
original coordinates: the positive diagonal entries tend to
zero, their sum is the trace, and the norm is bounded by that sum.

Let $E_0$ be the full confluent jet matrix with columns
$1,H_c,\ldots,H_c^{d-1}$. If $G_n^{\rm word}$ denotes the
received proof's quotient Gram in $1,X,\ldots,X^{d-1}$, then
\[
 G_n=(E_0^{-1})^*G_n^{\rm word}E_0^{-1},\qquad
 \det G_n=\frac{D_{n+1}}{|\det E_0|^2 B_{n-d+1}}\longrightarrow0.
 \tag{GC8}\label{gc:determinant}
\]
The numerator and denominator are the actual source and relation
Hankel determinants of GC5, with the original measure and no
rescaling. Explicitly $D_{n+1}$ is the Gram determinant of
$1,X,\ldots,X^n$ and $B_{n-d+1}$ is the Gram determinant of
$\nu,X\nu,\ldots,X^{n-d}\nu$, with $B_0=1$.
The combined basis
$1,\ldots,X^{d-1},\nu,\ldots,X^{n-d}\nu$ has determinant one
relative to the monomial basis because $\nu$ is monic.
Subtracting its attained relation projections gives the Schur
complement $G_n^{\rm word}$ and leaves the relation block.
The determinant therefore factors as
$D_{n+1}=B_{n-d+1}\det G_n^{\rm word}$.
The same quotient vector has original jet coordinates $E_0v$
when its word coordinates are $v$, proving the displayed
coordinate change and retaining $|\det E_0|^2$.

\subsection{An explicit sequence retaining every prescribed jet}
The convergence above admits a specified polynomial sequence.
For $u\in A$ let $P_u$ be its unique degree-below-$d$ representative
under $E$. Set $K=K_c$, $M_0=\mu(\R)=\|F_*\|^2$, and use the
actual suprema on $[-K,K]$:
\[
 f_\varepsilon(x)=-\frac{P_u(x)\nu(x)}{\nu(x)^2+\varepsilon^2},
 \qquad
 L_\varepsilon=
 \frac{\|P_u'\|_\infty}{2\varepsilon}
       +\frac{\|P_u\|_\infty\|\nu'\|_\infty}{\varepsilon^2}.
 \tag{GC13}\label{gc:regularized-inverse}
\]
The derivative of $f_\varepsilon$ has modulus at most
$L_\varepsilon$: use
$|\nu|/(\nu^2+\varepsilon^2)\le1/(2\varepsilon)$ and
$|\varepsilon^2-\nu^2|/(\nu^2+\varepsilon^2)^2\le1/\varepsilon^2$.
For $n\ge1$, put $y=(x+K)/(2K)$ and define the explicit
degree-at-most-$n$ polynomial
\[
 q_{\varepsilon,n}(x)=\sum_{k=0}^n
 f_\varepsilon(-K+2Kk/n)\binom nk y^k(1-y)^{n-k},
 \qquad P_{\varepsilon,n}=P_u+\nu q_{\varepsilon,n}.
 \tag{GC14}\label{gc:constructive-sequence}
\]
These exact coefficients use the original word variable; the
actual source is $P_{\varepsilon,n}(\mathsf H_c)F_*$.
One has $E P_{\varepsilon,n}=u$, including every derivative order,
and $\deg P_{\varepsilon,n}\le d+n$.
The variance of the Bernstein sampling point on $[-K,K]$ is
at most $K^2/n$, so Cauchy--Schwarz and the derivative bound give
$\|q_{\varepsilon,n}-f_\varepsilon\|_\infty
\le K L_\varepsilon/\sqrt n$.
Since $P_u+\nu f_\varepsilon=P_u\varepsilon^2/
(\nu^2+\varepsilon^2)$, the triangle inequality yields the
complete attained upper estimate
\[
 \begin{split}
 \langle G_{d+n}u,u\rangle^{1/2}
 &\le\|P_{\varepsilon,n}\|_{L^2(\mu)}\\
 &\le\left(\int
       \frac{|P_u(x)|^2\varepsilon^4}
            {(\nu(x)^2+\varepsilon^2)^2}\,d\mu(x)\right)^{1/2}
       +\frac{\|\nu\|_\infty\sqrt{M_0}\,K L_\varepsilon}{\sqrt n}.
 \end{split}
 \tag{GC15}\label{gc:constructive-bound}
\]
Taking $\varepsilon=1/j$ and $n=j^6$ gives a predetermined
sequence with original jets exactly $u$ and source norm tending
to zero. The second term is bounded by a fixed packet-dependent
constant times $j^{-1}$; the integral tends to zero by dominated
convergence and the zero measure of $\nu^{-1}(0)$ proved above.
No unproved quantitative rate for that integral is inserted.
This supplies explicit polynomial representatives for all the
vertical graph limits used next.

\subsection{The completed two-observation space retains the entire packet}
Consider now the actual two-observation map
\[
 P\longmapsto(\Phi(P),S(EP)),\qquad
 \|P\|_{\rm graph}^2=\|P\|_{L^2(\mu)}^2+\|S(EP)\|_{L^2(dx)}^2.
 \tag{GC9}\label{gc:graph}
\]
Both terms use the original source observation $dx$; the second
uses its already defined primitive section. No new form on the
packet is selected. Under GC5 its graph is $(P,EP)$ inside
$\mathcal H\oplus A_G$, where $A_G$ is $A$ with Gram $G=S^*S$.
Its closure is the entire direct sum. To prove this, a vector
$(f,v)$ orthogonal to that graph is first tested on $P=\nu Q$.
GC6 forces $f=0$. The surjectivity of $E$ then forces $v=0$.
The double orthogonal complement proves the claimed density.

The closure of the original relation ideal in this completed
graph is exactly $\mathcal H\oplus0$, by GC6 and GC2. Therefore
the complete quotient, with its attained metric, is
\[
 \frac{\overline{\{(P,EP):P\in\C[X]\}}}
      {\overline{\{(\nu Q,0):Q\in\C[X]\}}}
       =\frac{\mathcal H\oplus A_G}{\mathcal H\oplus0}
       \xrightarrow[\ [(f,u)]\mapsto u\ ]{\ \cong\ }A_G.
 \tag{GC10}\label{gc:completed-quotient}
\]
At finite cutoff its Gram in the original jets is exactly
$G+G_n$, since the second observation is constant on the fibre
$EP=u$. Hence its limit is $G$, and its determinant tends to
$\det G>0$. The forgetful projection onto the first source
observation has kernel $0\oplus A$. This explicitly identifies
every class lost by that projection, including all nilpotent jets.

The extension of the original reflected Weil form is equally
explicit:
\[
 \mathcal W((f,u),(h,v))
   =R_Z(u,J_gv)
   =\sum_{\rho\in Z}m_\rho\overline{u(1-\bar\rho)}v(\rho).
 \tag{GC11}\label{gc:weil}
\]
On polynomial graph vectors this is the original full Weil
pairing, since GC1 makes all other zero terms vanish. It extends
continuously by the finite projection to $A_G$; this extension
has radical $\mathcal H\oplus\prod_\rho(t_\rho)$.
The exact map relating GC9's positive observation norm to GC11
is the second projection, followed by the residue map and
$J_g$. For $r_f$ reflection-fixed roots and $r_p$ two-element
reflection orbits, its finite nonzero signature is
$(r_f+r_p,r_p)$, with $d-r_f-2r_p$ zero jet directions.
This follows from the value block $m_\rho$ at a fixed root and
the block $\left(\begin{smallmatrix}0&m_\rho\\m_\rho&0\end{smallmatrix}\right)$
at each pair. Thus no identification with the positive graph
norm is used to assign a sign to this arithmetic form.

Finally multiplication by $X$ on the dense polynomial graph
extends to $(M_x,M_{H_c})$ on $\mathcal H\oplus A_G$.
Its bound is $\max(K_c,\|M_{H_c}\|_G)$, since the two factors
have these bounds and $E(XP)=H_c EP$. The first component is
self-adjoint. The second has the precise original defect
\[
 S^*E_c-E_c^*S=M_{H_c}^*G-GM_{H_c},\qquad
 E_c=\mathsf H_cS-SM_{H_c}=\Theta\Psi_c.
 \tag{GC12}\label{gc:skew-defect}
\]
Expanding both sides and using $\mathsf H_c^*=\mathsf H_c$
proves the equality. Every discrepancy survives as the given
theta cochain and the retained packet component.
These are linear and Hilbert-space receivers of the already
defined full-support maps; no algebra structure on all of
$L^2(\mu)$ or equality of arithmetic forms is asserted.

\subsection{Finite-cutoff errors in the retained coordinates}
The convergence above also controls the complete finite Gram,
without a choice of a unit-mass source. Let $e_1,\ldots,e_d$ be
the original ascending jet basis of $A$, and apply GC13--14 to
each $u=e_i$, with $\varepsilon=j^{-1}$ and $n=j^6$. Denote the
resulting polynomial by $P_{i,j}$. Define the actual matrix
\[
 T_j=\left(\int\overline{P_{i,j}(x)}P_{l,j}(x)\,d\mu(x)
                    \right)_{i,l=1}^d,\qquad
 \alpha_j=\|G^{-1/2}T_jG^{-1/2}\|.
 \tag{GC16}\label{gc:finite-error}
\]
All entries use the original measure GC5 and the explicit
coefficients GC14. The lift $u\mapsto\sum_i u_iP_{i,j}$ has
observation $u$, so minimization proves
\[
 0<G_n\preceq T_j\preceq\alpha_jG
 \quad(n\ge d+j^6),\qquad \alpha_j\longrightarrow0.
 \tag{GC17}\label{gc:finite-domination}
\]
Indeed GC15 sends each diagonal of $T_j$ to zero, and its positive
trace bounds its operator norm. The fixed positive matrix $G$
then proves the last limit. No rate uniform in a growing packet
is inferred from this fixed-packet statement.

Write $\widetilde G_n=G+G_n$ for the exact graph quotient. In a
fixed original subspace frame $I$ with $r$ independent columns,
put $G^I=I^*GI$ and $\widetilde G_n^I=I^*\widetilde G_nI$.
The complete restricted error satisfies
\[
 0\le \log\frac{\det\widetilde G_n^I}{\det G^I}
       \le r\log(1+\alpha_j)\quad(n\ge d+j^6).
 \tag{GC18}\label{gc:restriction-error}
\]
Congruence by $I$ preserves GC17; congruence by $(G^I)^{-1/2}$
then puts every eigenvalue between $1$ and $1+\alpha_j$.
Taking their product proves GC18. In particular, at the four
unchanged indices $q-1,q,2q-1,2q$, whenever $q-1\ge d+j^6$,
\[
\begin{split}
0\ \le{}&\log\det\widetilde G_{q-1}^I
 +\log\det\widetilde G_q^I
 -\log\det\widetilde G_{2q-1}^I
 -\log\det\widetilde G_{2q}^I\\
 \le{}&2r\log(1+\alpha_j).
\end{split}
 \tag{GC19}\label{gc:four-cutoff}
\]
The lower bound uses the decreasing original minima separately
in the pairs $(q-1,2q-1)$ and $(q,2q)$. The fixed $\log\det G^I$
cancels with signs $+,+,-,-$; GC18 bounds the two positive
remainders and the negative remainders are nonnegative before
their signs are applied. For $r=0$ the empty determinant is one
and the displayed return is zero.

The connection to the one-observation minimum is exact at
\emph{every} admissible cutoff, including before GC19 applies:
\[
\begin{split}
\log\det\widetilde G_n^I
 &=\log\det(I^*G_nI)\\
 &\quad+\log\det\left(
 I_r+(I^*G_nI)^{-1/2}G^I(I^*G_nI)^{-1/2}\right).
\end{split}
 \tag{GC20}\label{gc:native-correction}
\]
Factor out the positive $I^*G_nI$ to prove this formula. Applying
the four signed evaluations to both sides retains the whole
second determinant. Thus the graph-completion calculation does
not assign an asymptotic to the original one-observation return.
Nor does it replace the original conductor source metric
SCD41: SCD37--42 is the exact isometry and four-cutoff receiver
for that source. Here GC20 identifies the metric change of the
specific two-observation construction GC9, including its entire
finite-cutoff correction.

\subsection{The prime-action defect on every spectral direction}
The theta cochain in GC12 carries quantitative information.
Use the inner product conjugate-linear in the first variable.
For any original eigenvector $u$ of $M_{H_c}$, with eigenvalue
$z_\rho$, self-adjointness of $\mathsf H_c$ gives
\[
 \operatorname{Im}\langle Su,E_cu\rangle
       =-\operatorname{Im}(z_\rho)\|Su\|^2,
 \qquad
 \|E_cu\|\ge |\operatorname{Im}z_\rho|\,\|Su\|.
 \tag{GC21}\label{gc:nonreal-defect}
\]
The first identity follows by substituting
$E_cu=\mathsf H_cSu-z_\rho Su$; the expectation of the
self-adjoint operator is real. Cauchy--Schwarz proves the bound.
Such a nonzero eigenvector is explicit even at a multiple root:
take the local idempotent times $(H_c-z_\rho)^{m_\rho-1}$.
The nonzero local derivative from GC2 proves it is nonzero and
that multiplication by $H_c$ acts on it by $z_\rho$.
For a nonfixed reflected pair, the two selected values are
distinct conjugates, hence nonreal. Thus GC21 detects its
nonzero discrepancy in the actual theta source, with the
original primitive norm and without changing its eigenvector.

The complete critical-point jets also remain visible. At a
reflection-fixed root with $m_\rho\ge2$, $z_\rho$ is real. Let
$e_\rho$ be its local idempotent and set
$u_0=e_\rho(H_c-z_\rho)^{m_\rho-1}$,
$u_1=e_\rho(H_c-z_\rho)^{m_\rho-2}$. They satisfy
$(M_{H_c}-z_\rho)u_0=0$ and
$(M_{H_c}-z_\rho)u_1=u_0\ne0$. Direct substitution gives
\[
\begin{split}
 \langle Su_0,E_cu_1\rangle-\langle E_cu_0,Su_1\rangle
       &=-\|Su_0\|^2,\\
 \|Su_0\|\|E_cu_1\|+\|E_cu_0\|\|Su_1\|
       &\ge\|Su_0\|^2>0.
\end{split}
 \tag{GC22}\label{gc:jet-defect}
\]
The two occurrences of $\mathsf H_c-z_\rho$ cancel by
self-adjointness, leaving exactly the displayed term and sign.
The triangle inequality proves the second line. This calculates
the retained higher-jet discrepancy even when the value of the
prime word is real. The value pairing GC11 kills these higher
jets, whereas GC22 proves that their source-action discrepancy
is nonzero. The residue projection in GC11 and the theta
cochain in GC12 are the exact maps connecting these statements.

\subsection{The lost positive class gives an explicit form-domain defect}
Choose a nonzero packet $u$ with $w=T_Z(u,u)>0$. Such a choice
exists for every nonempty reflected packet: at a fixed root use
its constant value vector, and at a nonfixed pair use the two
constant values $(1,1)$, putting all other coordinates zero.
The corresponding values of $w$ are $m_\rho$ and $2m_\rho$.
Let $P_u$ be its degree-below-$d$ representative and consider
the actual word subspace
\[
 \mathcal D_u=\C P_u+\nu\C[X]\subset L^2(\mu),\qquad
 Q_u(aP_u+\nu Q)=w|a|^2.
 \tag{GC23}\label{gc:form-domain}
\]
The decomposition has unique scalar $a$, since applying $E$
gives $au$. GC11 shows that $Q_u$ is precisely the restriction
of the original whole Weil form to this source subspace.
Every other zero contribution vanishes on the word family.
The domain is dense by GC6 and $Q_u$ is nonnegative.

Nevertheless this nonnegative form is not closable in the
original $L^2(\mu)$ source norm. The explicit polynomials
$P_j=P_{1/j,j^6}$ from GC14 satisfy
\[
 P_j\longrightarrow0\text{ in }L^2(\mu),\qquad
 Q_u(P_j-P_l)=0,\qquad Q_u(P_j)=w>0.
 \tag{GC24}\label{gc:form-defect}
\]
They all have observation $u$, so their differences belong
to the relation ideal. They are therefore Cauchy for the
form norm $\|P\|_{L^2(\mu)}^2+Q_u(P)$. If a closed extension
in the original source Hilbert space existed, completeness
would give a form-norm limit. Its source component must be
zero, whereas form-norm convergence would force its squared
form value to be $w$. A quadratic form has value zero at zero,
giving the contradiction.
The completed graph GC10 supplies the precise retained object:
this sequence converges there to $(0,u)$, whose second component
and form value $w$ survive. This establishes why the source
projection alone cannot carry the packet form, even on a
positive direction. It uses actual theta relations and the
specified sequence; it makes no inference that the full
Weil criterion has either sign.

\begin{figure}[ht]
\centering
\begin{tikzcd}[column sep=large,row sep=large]
 \C[X]\arrow[r,"{P\mapsto(P,EP)}"]\arrow[d,"\Phi"']
   &\mathcal H\oplus A_G\arrow[d,"\operatorname{pr}_1"]
      \arrow[r,two heads,"\operatorname{pr}_2"]&A_G\\
 \overline{\Phi(\C[X])}\arrow[r,"\cong"']&\mathcal H
\end{tikzcd}
\caption{The full graph closure is the direct sum, GC9--10.
The original theta relations are dense in the first source
component, GC6. The second observation retains every full jet
with the original primitive Gram. GC11 evaluates the arithmetic
form through that same second projection, and GC12 retains the
source-operator discrepancy. The lower isometry is the original
Mellin observation GC5.}
\end{figure}
